% !TEX root = lectures.tex


\section{Lecture April 18}

\subsection{Beyond Matrix Khintchine [Bandeira-Boedchardjo-von Handel '21]}

Recall the Matrix Khintchine inequality we started the class with.

Denote $X = \sum_i g_i A_i$ where $A_1, \dots, A_n \in \R^{d\times d}$ symmetric
and $g_i$ are iid standard Gaussians. Then,
\[ \E{\norm{X}_2} \lesssim \sigma_0 \sqrt{\log d} \]
and $\sigma_0 = \norm{\sum_i A_i^2}^{1/2}$. We also showed this is tight
with some simple diagonal examples. But, in some cases
we can remove the log factor (the expectation is always at least $\sigma_0^2$).
If we force them to be very non-commutative, we can perhaps hope for a better bound.

Note that the distribution of the vector $\sum_i g_i A_i$
is jointly Gaussian. So, it has a well-defined covariance $d^2 \times d^2$ matrix. Let's give this a name.
\begin{definition}
    \[ \sigma_1^2 = \norm{\Cov{X}}_2 \]
\end{definition}
Then, 
\begin{theorem}[BBvH'21]
    Under the same hypotheses,
    \[ \E{\norm{X}_2} \lesssim O(\sigma_0 + \log^{3/4}(d) \sigma_0^{1/2} \sigma_1^{1/2}) \le O(\sigma_0 + \log^{3/2}(d) \sigma_1) \]
\end{theorem}
In fact, you can get the tight constant on $\sigma_0$ from the previous claim.

However, this bound is not always better. Let's try applying this new inequality
to Kikuchi matrices. Take level $\ell$ Kikuchi matrices, $L_{C_1}, \dots, L_{C_m}$
where $C_i$'s are $4$-uniform hyperedges. Here, the size of the matrices is $\binom{n}{\ell} \times \binom{n}{\ell}$
where $\ell$ is polynomial in $n$, so $m = poly(n)$ is exponentially smaller than the size of the matrices.
However, for a given row $S$ in $\sum_i a_i L_{C_i}$, it has correlation $1$ with any $T$ that it has an edge to.
So, the covariance matrix will have really high spectral norm ($\binom{n - a}{\ell - b}$ neighbors for each vertex).

It's meaningful when for a given entry, there are very few entries we're correlated with (to make $\sigma_1$ small)
and the size of the matrix is reasonable (smaller than or around the same as) number of the summands. The large matchings
also correlate the entries together way too much.

For Wigner matrices, that is, a random symmetric matrix, then the histogram of the spectrum
look like a semicircle, centered at $0$ with radius $2\sqrt{n}$.

\subsubsection{Discrepancy Minimization}
We describe the simplest discrepancy minimization problem.
Let $a_1, \dots, a_n \in \R^d$ and all $\norm{a_i}_{\infty} \le 1$.
The goal is to find $b_1, \dots, b_n \in \{\pm 1\}$ such 
that $\norm{\sum_i b_i a_i}_{\infty} \le O(\sqrt{n}) \max\{1, \sqrt{\log \frac{d}{n}}\}$.
Note: When you randomly choose the $b_i$, then $\norm{\sum_i b_i a_i}_{\infty} \le \sqrt{n} \sqrt{\log d}$
with probability $\Omega(1)$.

\begin{theorem}[Spencer, 80s]
    Goal is achievable. This is proved non-constructively via the entropy method.
\end{theorem}

We can try to generalize this to the matrix setting.
$A_1, \dots, A_n \in \mathbb{R}^{n\times n}$, $\norm{A_i}_2 \le 1$ for all $i$.
The goal is to find $b_1, \dots, b_n \in \{\pm 1\}$
such that $\norm{\sum_i b_i A_i}_2 \le O(\sqrt{n} \max\{\sqrt{\log d/n}, 1\})$.
By matrix Khintchine, if we choose $b_1, \dots, b_n$ uniformly at random, then with probability $\Omega(1)$,
it's at most $O(\sqrt{n \log d})$. The conjecture is that this goal is acheivable, but we do not know how to prove it.

Today: we will show a paper by [Bansal-Jiang-Meka], the goal is achievable, if, in addition, $\norm{A_i}_F^2 \le \frac{n}{\log^3 n}$ for all $i$.
For today, we will take $d = n$ and the maximum that the squared Frobenius norm becomes is $n$.

Concrete question: Suppose $A_i$'s are adjacency matrices of perfect matchings. The operator norm is clearly $1$,
but the Frobenius norm is large. Is the Matrix Spencer conjecture true for this setting?

\begin{lemma}[Partial Coloring, Rothvoss '14]
    Let $\epsilon < 1/60000$, $\delta = \frac{3}{2}\epsilon \log_2 \frac{1}{\epsilon}$.
    Let $H\subseteq \R^n$ be a subspace of dimension at least $\ge (1 - \delta) n$.
    Let $K\subseteq H$ be a symmetric convex set such that $\gamma_H(K) \ge e^{-\delta n}$. Let
    $x_0 \in (-1, 1)^n$. Then, there is a point
    $x \in (x_0 +K) \subseteq [-1, 1]^n$ and $|\{i \mid |x_i| = 1\}| \ge \frac{\epsilon n}{2}$.
\end{lemma}

\begin{lemma}
    There are $c, c' > 0$ such that given $A_1, \dots, A_n \in \R^{d \times d}$,
    symmetric such that $\norm{\sum_i A_i^2}_{op} \le \sigma^2$ and $\sum_i \norm{A_i}_F^2 \le nf^2$, given
    $x_0 \in (-1 , 1)^n$, there is an $x \in [-1, 1]^n$ such that
    \[ \norm{\sum_{i= 1}^{n} (x_i - x_0,i) A_i}_{op} \le c(\sigma + \log^{3/4}(d) \sqrt{\sigma f}) \]
    and $|\{i \mid |x_i| = 1\}| \ge c'n$.
\end{lemma}

Consider the convex body
$K = \{x \in \R^n : \norm{\sum_i x_i A_i}_{op} \le c_0 \sigma + \log^{3/4}(d) \sigma^{1/2} f^{1/2}\} \subseteq \R^n$
for some constant $c_0$ we'll set later.
To prove its measure is large, we can use BBvH concentration inequality.
But first, we need to bound the covariance matrix
\[ \Cov X = \sum_i \vec{A_i} \vec{A_i}^T \implies \tr(Cov(X)) \le nf^2\]
But we have no idea what the largest eigenvalues are. But, we do know the number of eigenvalues
of value at least $f^2/\delta =: \Delta^2$ is at most $\delta n$. This gives us a subspace
$H = \{y \in \mathbb{R}^n : \sum_i y_i \Tr(A_i V_j) = 0, \forall j \in [k] \}$
where $V_1, \dots, V_k \in \R^{d\times d}$ are the eigenvectors of $\Cov(X)$ with eigenvalues at least $\Delta^2$.

Now, let's prove the claim. Sample $g \sim \R^n$ by first sampling $y \in H$,
$r \in H^{\perp}$, $g = y + r$.
Now, let $Y = \sum_i y_i A_i$ and $R = \sum_i r_i A_i$,
we have that $X = Y + R$ and $\E{X^2} = \E{Y^2} + \E{R^2}$. Then,
$\sigma_0(Y) \le \sigma_0(X) \le \sigma_0$. Then, we know that
for all $\vec{V_j}^{T} \Cov{Y} \vec{V_j} = 0$, then for all $v \perp \{V_1, \dots, V_k\}$
then $v^T \Cov{Y} v \le \Delta^2$. So,
\[ \E{\norm{Y}_{op}} \le c(\sigma_0(Y) +\log^{3/4}(d) \sigma_0(Y)^{1/2} \sigma_1(Y)^{1/2}) \le c'(\sigma_0 + \log^{3/4}(d) \sqrt{\sigma_0 f}) \]
Define $K' = K \cap H$. Then, taking $c_0 = 2c'$, the probability we land in $K'$ is thus at least $1/2$ by Markov,
i.e. $\gamma_H(K') \ge 1/2$. Furthermore, by construction and our previous bound on the
number of bad eigenvalues, the dimension of $H$ is at least $(1 - \delta)n$.
Thus, by our lemma , there is $x \in (x_0 + K') \cap [-1, 1]^n$ such that $|\{i \mid x_i \in \{\pm 1\}\}| \ge \epsilon/2$,
proving the claim.

Recall in our original theorem, we had control over the Frobenius squared norm creating $f \le \sqrt{\frac{n}{\log^3 n}}$,
giving us the right answer by using the claim. Now, we can simply recurse the lemma over and over (by
restricting the coordinates to those which are not set), getting a partial coloring each time.
Furthermore, since we are setting $\Omega(n)$ coordinates each time, the error bound we get is $\sqrt{\text{number of coordinates left}}$,
thus the total sum of operator residuals is something like $\sqrt{n} + \sqrt{cn} + \sqrt{c^2 n} + \dots = O(\sqrt{n})$.

\subsubsection{A Proof of the Partial Coloring Lemma}
Finally, let's prove Rothvoss' result.
\begin{theorem}
    Let $\epsilon, \delta$ be the same as in the Partial coloring lemma. Let $K\subseteq \R^n$
    by a symmetric convex body, $\gamma_n(K) \ge e^{-\delta n}$.
    Let $x* \sim N(0, 1)^n$ and $y^* \in K\cap [-1, 1]^n$ minimizing $\norm{x* - y*}_2$.
    Then, with probability $ 1- e^{-\Omega(n)}$, $y*$ has $\ge \frac{\epsilon n}{2}$ coordinates set to $\pm 1$.
    (Notice this gives an algorithm to find the point, as long as $K$ has a polynomial time projection
    oracle).
\end{theorem}

\begin{fact}[Sidak, Khatri '67, Special case of Gaussian Correlation Inequality]
    If $K \subseteq \R^n$ symmetric, convex in $\R^n$ and let $S\subseteq \R^n := \{x \in \R^n \mid |\langle v, x \rangle| \le \lambda\}$
    be a strip. Then, $\gamma_n(K\cap S) \ge \gamma_n(K) \gamma_n(S)$.
\end{fact}

\begin{fact}
    Let $h(\epsilon) = \epsilon \log_2 \frac{1}{\epsilon} + (1 - \epsilon) \log_2 \frac{1}{1 - \epsilon}$.
    Then for all $0 < \epsilon <1/2$,
    then the number of sets $|I| \le \epsilon n$ is at most $2^{nh(\epsilon)} \le 2^{3\epsilon/2 \log_2 1\epsilon}$.
\end{fact}

\begin{fact}[Gaussian Isoperimetric Inequality]
    $K \subseteq \R^n$, $H$ halfspace with $\gamma_n(K) = \gamma_n(H)$.
    Then, for all $\delta > 0$, $\gamma_n(K_{\delta}) \ge \gamma_n(H_{\delta})$
    where $S_{\delta}$ is the $\delta$-net of $K$, i.e. the set of all points that are distance at most $\delta$ from $S$.
\end{fact}

\begin{corollary}
    Let $\epsilon > 0$. Then for all $K$, $\gamma_n(K) \ge e^{-\epsilon n}$,
    then $\gamma_n(K_{3\sqrt{\epsilon n}}) \ge 1 - e^{-\epsilon n}$.
\end{corollary}

To prove this corollary, just prove it for halfspaces! To get this volume, my halfspace must be distance $3/2 \sqrt{\epsilon n}$
from the center, so if we expand to $3\sqrt{\epsilon n}$, we get most of the volume. Now we can prove the theorem.
The intuition is that we can drop non-tight constraints from the solution $y$, and so the volume of the body will be much larger than
the volume of $[-1, 1]^n$.

\begin{proof}
    Claim 1: $d(x*, [-1, 1]^n) \ge c\sqrt{n}$ with probability $1 - e^{-\Omega(n)}$.
    Let $I^*$ be the coordinate constraints which are tight (where $y^* = \pm 1$).
    Projecting to our set is the same as projecting to the set $K(I*) = K \cap \{ x \in \R^n: |x_i| \le 1, \forall i \in I*\}$.
    Suppose $|I*| <\epsilon n$. By Sidak/Khatri, 
    \[ \gamma_n(K(I*)) \ge \gamma_n(K) \prod_{i \in I*} \gamma_n(\text{strip}) \ge e^{-\delta n}
     e^{-\epsilon n/2} \ge e^{-2\delta n} \]
    Union-bounding over all $I*$ we get $\gamma_n(\cup_{I*} K(I*)) \le e^{-\delta n}$.
    Thus, $\gamma_n(K(I*)_{3\sqrt{2\delta n}}) \ge 1 -e^{-\delta n}$ and we can set $\delta$ small enough.
\end{proof}